\solutionsheader

\solution{Черный пакет}{15.0}

Все измерения будем проводить с помощью батареи, то есть заряжая ею черный ящик. Поэтому измерим его ЭДС: $V_0 = \qty{8.89}{\V}$. 

Далее подключим черный ящик последовательно к батарее и амперметру. Измерим силу тока через черный ящик: $I_r = \qty{1.81}{\mA}$. Попробуем поменять полярность батареи и измерить силу тока в этом случае. Дождавшись момента, когда ток устновится, видим, что сила тока иммет значение равное: $I_l = \qty{0.06}{\mA}$. Факт того, что сила тока разная в разных направлениях указывает на то, что в цепи имеется диод. Данное значениетока  очень близко к нулю и заметно мало по сравнению. Поэтому можем считать, что при такой полярности ток в цепи не течет и диод соединен последовательно с резистором. Таким образом, одна ветвь черного ящика это диод и резистор соединенные последовательно. При этом, резисторов, соединенных к выводам нет, так как во втором случае ток был нулевым.

Из значения тока можем с легкостью найти сопротивление первого резистора: 
\begin{eq}
    R_1 = \frac{V_0}{I_r} \approx \qty{4.91}{\kOhm}
\end{eq}

Во время подключения амперметра, можем заметить, что ток менялся со временем и приходил к стационарному значению лишь через некоторое время. Такое поведение получается при присутствии конденсатора. Кроме этого, из показании тока в двух направлениях, можно сделать вывод, что к выводам черного ящика соединены либо резисторы с диодом или конденсаторы, последовательно соединенные с другими элементами. Так как несколько параллельно соединенных ветвей одинакового типа будут себя вести схоже, они будут эквивалентны самому упрощенной схеме, где имеется только две ветви: одна резистор и диод, другая конденсатор с другими элементами.  

Для изучения поведения цепи в зависимости конденсатора поступим следующим образом: Подключим  параллельно черный ящик, батарею и вольтметр. Дождемся некоторое время(время такого порядка, за которое у нас при подключении амперметра устанавливался ток) и отсоединим батарею от цепи. После будем исследовать зависимость напряжения от времени. 

Нарисуем схему черного ящика, насколько нам известно. Также нарисуем вольтметр, чтобы понять, что мы мерим: 

\cpic{0.4}{figures/test.jpg}

Из схемы становится яснее то, что измеряемое нами напряжение это падения напряжения на резисторе. 

\begin{table}[h!]
\centering
\caption{Данные разряда (часть 1)}
\resizebox{\textwidth}{!}{
\begin{tabular}{l|*{16}{c}}
\hline
$V$, В & 1.263 & 1.096 & 0.923 & 0.861 & 0.548 & 0.490 & 0.448 & 0.390 &0.371& 0.354 & 0.338 & 0.326 & 0.318 & 0.304 & 0.296 & 0.286 \\
$t$, с & 2.07 & 2.39 & 2.94 & 3.39 & 5.23 & 5.96 & 6.97 & 8.45 &9.49& 10.47 & 11.47 & 12.81 & 14.09 & 15.61 & 17.26 & 18.37 \\
$\ln V$ & 0.233 & 0.092 & -0.080 & -0.150 & -0.601 & -0.713 & -0.803 & -0.942 &-0.992&  -1.038 & -1.085 & -1.121 & -1.146 & -1.191 & -1.217 & -1.252 \\
\hline
\end{tabular}
}
\end{table}

Нарисуем график данного зависимости:

\begin{table}[h!]
\centering
\caption{Данные разряда (часть 2)}
\resizebox{\textwidth}{!}{
\begin{tabular}{l|*{14}{c}}
\hline
$V$, В & 0.279 & 0.265 & 0.258 & 0.246 & 0.234 & 0.227 & 0.220 & 0.215 & 0.210 & 0.206 & 0.197 & 0.192 & 0.187 & 0.179 \\
$t$, с & 20.47 & 23.03 & 26.58 & 30.23 & 36.41 & 41.07 & 45.62 & 50.65 & 54.99 & 59.57 & 65.20 & 75.23 & 82.92 & 108.51 \\
$\ln V$ & -1.277 & -1.328 & -1.355 & -1.402 & -1.452 & -1.483 & -1.514 & -1.537 & -1.561 & -1.580 & -1.625 & -1.650 & -1.677 & -1.720 \\
\hline
\end{tabular}
}
\end{table}

\vspace{1cm}

\section*{Графики}

\begin{figure}[h!]
    \centering
    % График V(t)
    \begin{tikzpicture}
    \begin{axis}[
        title={Зависимость напряжения от времени $V(t)$},
        xlabel={Время $t$, с},
        ylabel={Напряжение $V$, В},
        grid=major,
        width=12cm,
        height=8cm,
        mark size=1.5pt
    ]
    \addplot[
        only marks,
        mark=*,
        color=blue
    ] coordinates {
        (2.07, 1.263) (2.39, 1.096) (2.94, 0.923) (3.39, 0.861) (5.23, 0.548) (5.96, 0.490) (6.97, 0.448) (8.45, 0.390) (10.47, 0.354) (11.47, 0.338) (12.81, 0.326) (14.09, 0.318) (15.61, 0.304) (17.26, 0.296) (18.37, 0.286) %(20.47, 0.279) (23.03, 0.265) (26.58, 0.258) (30.23, 0.246) (36.41, 0.234) (41.07, 0.227) (45.62, 0.220) (50.65, 0.215) (54.99, 0.210) (59.57, 0.206) (65.20, 0.197) (75.23, 0.192) (82.92, 0.187) (108.51, 0.179)
    };
    \end{axis}
    \end{tikzpicture}
\end{figure}


\begin{figure}[h!]
    \centering
    % График V(t)
    \begin{tikzpicture}
    \begin{axis}[
        title={Зависимость напряжения от времени $V(t)$},
        xlabel={Время $t$, с},
        ylabel={Напряжение $V$, В},
        grid=major,
        width=12cm,
        height=8cm,
        mark size=1.5pt
    ]
    \addplot[
        only marks,
        mark=*,
        color=blue
    ] coordinates {
        (20.47, 0.279) (23.03, 0.265) (26.58, 0.258) (30.23, 0.246) (36.41, 0.234) (41.07, 0.227) (45.62, 0.220) (50.65, 0.215) (54.99, 0.210) (59.57, 0.206) (65.20, 0.197) (75.23, 0.192) (82.92, 0.187) (108.51, 0.179)
    };
    \end{axis}
    \end{tikzpicture}
\end{figure}

Графики напряженности довольно похожи на экспоненциальную, особенно первый график, который показывает начало процесса. Поэтому будем рассматривать этот участок. Нарисуем график для этих точек в координатах $\ln V, t$. 

\begin{figure}[h!]
    \centering
    \begin{tikzpicture}
    \begin{axis}[
        title={Полулогарифмический график $\ln V$ от $t$},
        xlabel={Время $t$, с},
        ylabel={$\ln V$},
        grid=major,
        width=12cm,
        height=8cm,
        mark size=1.5pt
    ]

    % 2. Линия тренда для начального участка (2-6 с)
    \addplot[
        domain=1.5:6.5,
        samples=2,
        color=red,
        thick,
        dashed
    ] {-0.238*x + 0.668};
    \addlegendentry{Аппроксимация}
    
    \addplot[
        only marks,
        mark=triangle*,
        color=purple
    ] coordinates {
        (2.07, 0.233) (2.39, 0.092) (2.94, -0.080) (3.39, -0.150) (5.23, -0.601) (5.96, -0.713) (6.97, -0.803) (8.45, -0.942)(10.47, -1.038) (11.47, -1.085) (12.81, -1.121) %(14.09, -1.146) (15.61, -1.191) (17.26, -1.217) (18.37, -1.252) (20.47, -1.277) (23.03, -1.328) (26.58, -1.355) (30.23, -1.402) (36.41, -1.452) (41.07, -1.483) (45.62, -1.514) (50.65, -1.537) (54.99, -1.561) (59.57, -1.580) (65.20, -1.625) (75.23, -1.650) (82.92, -1.677) (108.51, -1.720)
    };
    \end{axis}
    \end{tikzpicture}
\end{figure}

График оказался не линейным. Однако, мы имеем линейные участки, особенно в начале процесса. Последующее отклонение от линейности считается вызванным из за неидеальности диода.  

Аппроксимируем прямую область линейной функцией $\ln V = -kt + b$. Найдем параметры функции по МНК: $k \approx \qty{-0.238}{\per\s}; b = 0.668 $. 

Характерное время, $\tau $ равняется следущему выражению: $\tau = C\cdot (R_1 + R_x)$, где $R_x$ сопротивление неизвестной части $X$. На самом деле, на поведении черного ящика сказывается только сопротивление элемента $X$. Все другие схемы будут эквивалентны, потому что присутствие остальных элементов не скажется на поведении. Таким образом окончательная схема выглядит следующим образом: 

\cpic{0.4}{figures/black_box.jpg}

В этой сборке зависимость напряжения на выходе имеет вид: 
\begin{eq}
    V(t) = V_0 \frac{R_1}{R_1+R_2}\cdot \exp{(-t/\tau)} \implies \ln V = \ln V_0 \frac{R_1}{R_1+R_2} - \frac t \tau.
\end{eq}
Найдем сперва сопротивление второго резистора: 
\begin{eq}
    R_1 \cdot \left(\frac{V_0}{e^b} - 1\right) \approx\qty{17.4}{\kOhm}
\end{eq}

Теперь, имея значения сопротивления, можем найти значения ёмкости конденсатора: $C = \qty{188}{\uF}$. 